\chapter*{Conclusion générale}
\addcontentsline{toc}{chapter}{Conclusion générale}

Dans le cadre de notre deuxième projet de fin d’année à l’École Nationale Supérieure d’Ingénieurs de Tunis (ENSIT), nous avons conçu et développé une plateforme intelligente de pépinière en ligne nommée \textbf{PhytoCare}. Cette plateforme combine une solution e-commerce dédiée à la vente de plantes avec un système intelligent de diagnostic et de recommandation basé sur l’intelligence artificielle.

Ce projet nous a permis de mettre en pratique les connaissances acquises durant notre formation, notamment en analyse et conception des systèmes, en développement web, en intelligence artificielle ainsi qu’en apprentissage automatique. Il nous a également permis de mieux comprendre les différentes étapes de réalisation d’un projet informatique, depuis l’étude des besoins jusqu’à la phase de conception, d’implémentation et de validation.

À travers cette application, nous avons proposé une solution permettant aux utilisateurs de consulter et d’acheter des plantes en ligne, tout en bénéficiant d’un système intelligent capable d’analyser l’état de santé des plantes à partir d’images et de fournir des recommandations adaptées.

Ce projet a constitué une expérience enrichissante aussi bien sur le plan technique que sur le plan humain. Il nous a permis de développer notre esprit d’équipe, notre autonomie ainsi que notre capacité à résoudre des problèmes concrets.

Bien que les objectifs principaux du projet aient été atteints, plusieurs améliorations peuvent encore être envisagées afin d’enrichir davantage les fonctionnalités de la plateforme et d’améliorer l’expérience utilisateur. Parmi les perspectives futures, nous envisageons l’amélioration des performances du modèle de diagnostic intelligent à l’aide de datasets plus riches et plus variés, l’intégration d’un système de suivi personnalisé des plantes pour les utilisateurs, ainsi que l’ajout d’un système de paiement en ligne sécurisé.

Nous proposons également le développement d’une application mobile afin de faciliter l’accès à la plateforme et d’offrir une meilleure accessibilité aux utilisateurs. Enfin, l’intégration d’un système de recommandations plus avancé, basé sur les préférences et l’historique des utilisateurs, permettrait de rendre la plateforme plus intelligente et plus adaptée aux besoins de chaque utilisateur.

En conclusion, ce projet représente une étape importante dans notre parcours académique et constitue une expérience formatrice aussi bien sur le plan technique que sur le plan humain. Il nous a permis de consolider nos connaissances en développement web, en intelligence artificielle ainsi qu’en apprentissage automatique, tout en développant notre capacité à travailler en équipe et à gérer les différentes phases d’un projet informatique. Cette expérience sera sans aucun doute très bénéfique pour la réalisation de nos futurs projets académiques et professionnels.